\section{A mixed trace that detects the valuation}\label{sec:mixedtrace}
The valuation coordinate can be read from a joint word in the same representation. For $T(g)=\tr_{\Oloc}\rho(g)$, define
\[
D(\tau,g)=T(\tau g)-T(\tau)-T(g)+2.
\]
The elementary identity
\begin{equation}\label{eq:mixedidentity}
D(\tau,g)=\tr\bigl((\rho(\tau)-I)(\rho(g)-I)\bigr)
\end{equation}
shows why the product label matters. Two lists of individual traces do not determine the trace of the product.

The local Kummer basis $[\phi],[3],[7]$ supplies a unit annihilating probe $\tau_j$ after adjoining the required $7^j$ roots of unity. It kills the unit coordinates and evaluates the uniformizer coordinate to one. The inherited construction uses the cyclotomic action to show restriction remains injective. At higher cyclotomic levels this is not a blanket prime to seven degree argument, because the extension also has a seven primary part.

To define the valuation functional used here, strictly align the actual local toric cyclotomic line at $K_7$ with the reference line. The aligned local representative of $X$ has $p=\beta=0$ and upper entry $q$ a $k(\chi)$ Kummer cocycle. Define $\nu_7(X)$ to be its $[7]$ valuation coordinate. Remaining aligned strict gauges change $q$ by a coboundary and a scalar multiple of residual $\alpha$, both of valuation zero. Thus the coordinate is well defined under these strict equivalences. The reference $b_0$ has valuation zero and the actual difference is $b_7-b_0=7q$, giving the fixed value below. This local alignment never erases beta from the global deformation.

The actual period valuation matrix is $7cI$, so this probe has
\[
\rho(\tau_j)=I+7cPE_{12}P^{-1}\pmod{7^j},
\]
where the same local toric line has the required alignment with the natural line. At actual depth one and a detecting $g$ in the residual kernel, \eqref{eq:mixedidentity} gives
\begin{equation}\label{eq:mixedvalue}
D(\tau,g)=49\nu_7(X)\beta(g)\pmod{343}.
\end{equation}
The actual fixed seven period comparison requires
\[
\nu_7(X)=c\frac{v}{2\epsilon}\quad\text{in }\F_7.
\]
If $\beta(g)\ne0$, the mixed trace determines this value.

The all depth formula uses $t=1+v_7(c)$, precision $j\ge n+t+1$, the probe $\tau_j$, and a detecting $g\in\ker\overline\rho$ with $\beta(g)\ne0$. It is
\[
v_7(D(\tau_j,g))=n+t,
\qquad
\frac{D(\tau_j,g)}{7^{n+t}}
=\frac{c}{7^{t-1}}\frac{v}{2\epsilon}\beta(g)\pmod7.
\]
The coefficient on the right is nonzero in the detecting case. This is an exact valuation statement under the probe and actual geometric hypotheses. It does not bound $n$ or $c$.

\begin{example}
Let $U_a=\left(\begin{smallmatrix}1&7a\\0&1\end{smallmatrix}\right)$ and $V=\left(\begin{smallmatrix}1&0\\7&1\end{smallmatrix}\right)$. Both have trace two and determinant one, for every $a$. Their product has trace $2+49a$. Thus $a=1,2$ give product traces $51,100$ modulo $343$, while every individual trace and every trace modulo forty nine loses this distinction.
\end{example}
These matrices are exact controls for the information loss. They are not claimed to be the actual labelled probe and detecting element. The current eight tuples do not supply $T(\tau)$, $T(g)$, and $T(\tau g)$ from one numerical global realization. The required mixed value is a target, not an observed contradiction.

For instance, $c=1,v=3,\epsilon=-1,\beta(g)=1$ requires $D\equiv98\pmod{343}$. Writing that number does not evaluate the trace of an unknown candidate. A future global calculation must identify the words and the representation together.
